\chapter{Percutaneous Coronary Intervention}

\section*{Introduction \& Clinical Indications}
Percutaneous coronary intervention (PCI) uses a catheter, balloon, and usually a coronary stent to restore flow through a narrowed or blocked coronary artery. It is central to the treatment of selected acute coronary syndromes and can improve symptoms in appropriately selected patients with stable coronary disease. The decision between PCI, coronary artery bypass grafting, and medical therapy depends on anatomy, urgency, ischemic burden, diabetes, ventricular function, kidney function, bleeding risk, and patient preference.

Primary PCI is an emergency reperfusion strategy for eligible patients with acute ST-elevation myocardial infarction when timely treatment is available. In other settings, angiography and a multidisciplinary heart-team discussion may be needed before revascularization. PCI is not a cure for atherosclerosis; secondary prevention remains essential.

\section*{Relevant Surgical Anatomy}
The left and right coronary arteries arise from the aortic root and divide into vessels supplying the myocardium. Lesion location, calcification, vessel diameter, bifurcation anatomy, chronic total occlusion, and collateral circulation affect technical difficulty and treatment choice. The radial and femoral arteries provide access routes, while the coronary ostia and aortic root must be engaged safely with a guiding catheter.

\section*{Preoperative Workup \& Preparation}
Assess symptoms, electrocardiography, hemodynamics, renal function, blood count, bleeding risk, allergies, and antiplatelet or anticoagulant use. Coronary angiography defines anatomy; intravascular ultrasound or optical coherence tomography may assist with lesion assessment and stent optimization. Explain bleeding, vascular injury, contrast-associated kidney injury, arrhythmia, stroke, emergency surgery, stent thrombosis, restenosis, myocardial infarction, and death. Antiplatelet treatment is planned before and after PCI according to the indication and stent strategy.

\section*{Surgical Technique \& Procedure}
After local anesthesia and monitoring, arterial access is obtained, most commonly through the radial artery when suitable. A guiding catheter is positioned at the coronary origin, a wire crosses the lesion, and balloon angioplasty prepares the segment. A drug-eluting stent is deployed when indicated and expanded to restore vessel diameter. Calcified or complex lesions may require specialized plaque-modification techniques. Final angiography assesses flow, dissection, perforation, and residual narrowing before sheath removal and hemostasis.

\section*{Complications \& Risks}
Complications include access-site bleeding or hematoma, radial occlusion, femoral pseudoaneurysm, coronary dissection, perforation, no-reflow, acute stent thrombosis, contrast reaction, kidney injury, stroke, arrhythmia, myocardial infarction, and emergency need for surgery. Interruption of prescribed antiplatelet therapy can increase stent-thrombosis risk and should be discussed with the treating team before any procedure.

\section*{Postoperative Care \& Recovery}
Monitor the access site, distal circulation, rhythm, symptoms, kidney function, and signs of recurrent ischemia. Radial or femoral precautions depend on the access method. Prescribed antiplatelet therapy, lipid lowering, blood-pressure and diabetes management, smoking cessation, exercise, and cardiac rehabilitation reduce future risk. New chest pain, shortness of breath, syncope, or uncontrolled bleeding requires urgent assessment.

\section*{Outcomes \& Prognosis}
PCI can rapidly restore coronary flow and relieve angina, but long-term outcome depends on infarct size, ventricular function, disease complexity, adherence to prevention, and the completeness of revascularization. CABG may offer better outcomes for selected left-main, multivessel, or diabetic disease; treatment should be individualized through a heart-team approach.

\section*{References}
\begin{enumerate}
\item American College of Cardiology/American Heart Association/SCAI. Guideline and consensus documents on coronary revascularization and PCI.
\item American Heart Association. 2025 Guideline for the Management of Patients With Acute Coronary Syndromes.
\item European Society of Cardiology. Guidelines for acute coronary syndromes and chronic coronary syndromes.
\end{enumerate}
